%%%%%%% ESTRUCTURA DEL TFG %%%%%%%%%%%%%%
% A.PORTADA (según la estructura indicada a continuación) 
% B.CONTRAPORTADA (según la estructura indicada a continuación) 
% C.ÍNDICE GENERAL DE CONTENIDOS  
% D.ÍNDICE DE TABLAS  
% E.ÍNDICE DE FIGURAS 
% F.RESUMEN (Podrá incluirse también en inglés, si así lo indica el Tutor1) 
% G.CUERPO DEL TRABAJO (según la estructura indicada a continuación) 
% H.REFERENCIAS BIBLIOGRÁFICAS 
% (Según norma ISO690)
% I.ANEXOS, si los hubiera
%%%%%%%%%%%%%%%%%%%%%%%%%%%%%%%%%%%%%%%
%%%%%% Cuerpo del trabajo
% 1.INTRODUCCIÓN 
% 2.OBJETIVOS 
% 3.ANTECEDENTES / ESTADO DEL ARTE 
% 4.MÉTODOLOGÍA 
% 5.IMPLEMENTACIÓN Y DESARROLLO (Cuando proceda) 
% 6.RESULTADOS Y DISCUSIÓN 
% 7.CONCLUSIONES 
%%%%%%%%%%%%%%%%%%%%%%%%%%%%%%%%%%%%%%%%%%%
%%%%%%%%%%%%%%%%%%%%%%%%%%%%%%%%%%%%%%%%%%%
% Tamaño: Normalizado UNE A-4, salvo planos. 
% Tipo y tamaño de letra del texto: Times New Roman 12 pt, Arial 12 pt o similar. 
% Interlineado del texto: 1,5 líneas. 
% Márgenes del texto: Superior, Inferior y Derecha, 2,5 cm; Izquierda, 4 cm. 
% Numeración de páginas en margen inferior derecha y tamaño 8 pt. 
% Las  figuras  serán  numeradas  y  tituladas  debajo  de  las  mismas  (indicando  su  fuente  si  no  son  de  elaboración propia). 
% Las  tablas  serán  numeradas  y  tituladas  encima  de  las  mismas  (indicando  su  fuente  si  no  son  de  elaboración propia). 
%%%%%%%%%%%%%%%%%%%%%%%%%%%%%%%%%%%%%%%%%%%
%%%%%%%%%%%%%%%%%%%%%%%%%%%%%%%%%%%%%%%%%%%
%%%%%%%%%%%%%%%%%%%%%%%%%%%%%%%%%%%%%%%%%%%
%%%%%%%%%%%%%%%%%%%%%%%%%%%%%%%%%%%%%%%%%%%

%\documentclass[a4paper]{article}
\documentclass[12pt,a4paper,twoside]{book} % tamaño 12
%\usepackage[spanish]{babel}
\usepackage[spanish, es-tabla]{babel}
\usepackage[utf8x]{inputenc}
\usepackage[table]{xcolor} % para poner color en las tablas
\usepackage{tikz} % grafos
\usetikzlibrary{arrows,shapes,positioning,shadows,trees}
\usepackage{amsmath} % códigos
\usepackage{setspace} %interlineado algoritmo
\usepackage{algorithm}
\usepackage{algorithmicx}
\usepackage{algpseudocode}
\algdef{SE}[DOWHILE]{Do}{doWhile}{\algorithmicdo}[1]{\algorithmicwhile\ #1}%
\usepackage{subfig} % Para poner imágenes
\usepackage{graphicx} % para poner dos imágenes juntas
\usepackage[colorinlistoftodos]{todonotes}
\usepackage{anysize} 
\marginsize{4cm}{2.5cm}{2.5cm}{2.5cm} 
% Márgenes {izquierda}{derecha}{arriba}{abajo}. 
\usepackage{mathptmx}% http://ctan.org/pkg/mathptmx Times new roman
\usepackage{amsfonts}% Para expresiones matemáticas
%\usepackage{setspace} % interlineado
\renewcommand{\baselinestretch}{1.5} % Interlineado
\usepackage{enumerate}% para los enumerate
\usepackage{fancyhdr}
\pagestyle{fancy}
\usepackage[none]{hyphenat} %saltos y justificación texto
\usepackage{float} 
\usepackage{tikz}
\usetikzlibrary{calc}
\usepackage{changepage}
\usepackage{ragged2e}

%%%%%%%%%%%%%%%%%%%%%%%
%%%%%%%% XML %%%%%%%%%%
\usepackage{listings}

\usepackage{appendix}
% cambiar apéndices por anexos
\renewcommand{\appendixname}{Anexos}
\renewcommand{\appendixtocname}{Anexos}
\renewcommand{\appendixpagename}{Anexos}

\usepackage{color}
\definecolor{gray}{rgb}{0.4,0.4,0.4}
\definecolor{darkblue}{rgb}{0.0,0.0,0.6}
\definecolor{cyan}{rgb}{0.0,0.6,0.6}

\lstset{
  basicstyle=\ttfamily,
  columns=fullflexible,
  showstringspaces=false,
  commentstyle=\color{gray}\upshape
}

\lstdefinelanguage{XML}
{
  morestring=[b]",
  morestring=[s]{>}{<},
  morecomment=[s]{<?}{?>},
  stringstyle=\color{black},
  identifierstyle=\color{darkblue},
  keywordstyle=\color{cyan},
  morekeywords={id,tx,ty,tz,rx,ry,rz,scale,file,type}% list your attributes here
}
%%%%%%%%%%%%%%%%%%%%%%%%%%%%
%%%%%%%%%%%%% COLOR %%%%%%%%
\definecolor{palegreen}{rgb}{0.6, 0.98, 0.6}
\definecolor{palegoldenrod}{rgb}{0.93, 0.91, 0.67}

\makeatletter
\def\BState{\State\hskip-\ALG@thistlm}
\makeatother


\rhead{\begin{picture}(0,0) \put(0,0){\includegraphics[width=1cm]{pictures/logouex_transp.png}}
\end{picture}}


\begin{document}
\sloppy 
\pagebreak 
%http://tex.stackexchange.com/questions/129088/create-a-frame-for-a-title-page
\makeatletter

\begin{titlepage}
	\begin{adjustwidth}{-1.5cm}{0cm}
		\begin{tikzpicture}[remember picture, overlay]
		\usetikzlibrary{calc}
		\draw[line width = 1.5pt] ($(current page.north west) + (1in,-1in)$) rectangle ($(current page.south east) + (-1in,1in)$);
		\end{tikzpicture}
		
		\vspace{-3.5em}
		\hspace{-0.5em}
		\begin{minipage}{0.45\textwidth}
			\begin{flushleft}
				\includegraphics[width=1.75cm]{pictures/logouex_transp.png}
			\end{flushleft}
		\end{minipage}
		
		\vspace{-6em}
		\hspace{20em}
		\begin{minipage}{0.45\textwidth}
			\begin{flushright}
				\includegraphics[width=4cm]{pictures/logoEpcc.png}
			\end{flushright}
		\end{minipage}\\[1.5cm]
		
		\begin{center}
			
			
			% Upper part of the page
			\textsc{\LARGE UNIVERSIDAD DE EXTREMADURA}\\[4cm]
			
			\textsc{\Large Escuela Politécnica}\\[0.5cm]
			\textsc{\Large Nombre de los Estudios}\\[3cm]
			\textsc{\Large Trabajo Fin de Grado}\\[0.5cm]
			{ \large \bfseries Título trabajo fin de grado}\\[4.0cm]
			\vfill
			% Bottom of the page
			{\large}
		\end{center}
	\end{adjustwidth}
\end{titlepage}


\pagestyle{empty}


\begin{titlepage}
	\begin{adjustwidth}{-1.5cm}{0cm}
		\begin{tikzpicture}[remember picture, overlay]
			\usetikzlibrary{calc}
			\draw[line width = 1.5pt] ($(current page.north west) + (1in,-1in)$) rectangle ($(current page.south east) + (-1in,1in)$);
		\end{tikzpicture}

		\vspace{-2em}
		\hspace{-0em}
		\begin{minipage}{0.45\textwidth}
			\begin{flushleft}
				\includegraphics[width=1.60cm]{pictures/logouex.jpg}
			\end{flushleft}
		\end{minipage}

		\vspace{-6.5em}
		\hspace{19.5em}
		\begin{minipage}{0.45\textwidth}
			\begin{flushright}
				\includegraphics[width=3cm]{pictures/logoEpcc.jpg}
			\end{flushright}
		\end{minipage}\\[1.5cm]

		\begin{center}


			% Upper part of the page
			\textsc{\LARGE UNIVERSIDAD DE EXTREMADURA}\\[2cm]

			\textsc{\Large Escuela Politécnica}\\[0.5cm]
			\textsc{\Large Grado en Ingeniería Informática}\\[0.5cm]
			\textsc{\Large en Ingeniería del Software}\\[2.5cm]
			\textsc{\Large Trabajo Fin de Grado}\\[0.5cm]
			{ \large \bfseries Lenguaje especifico de modelado}\\[0.5cm]
			{ \large \bfseries de sistemas Blockchain sobre TrustOS}\\[2.0cm]
			{ \large Autor}\\[0.5cm]
			{ \large \bfseries Pedro Miguel Carmona Broncano}\\[0.5cm]
			{ \large  Tutor }\\[0.5cm]
			{ \large \bfseries Pedro José Clemente Martín}\\[0.5cm]
			\vfill
			% Bottom of the page
			{\large}
		\end{center}
	\end{adjustwidth}
\end{titlepage}
\newpage

\makeatother
\thispagestyle{empty}
\newpage
\thispagestyle{empty}


\tableofcontents
\thispagestyle{chapterpage}
\newpage
\listoftables
\thispagestyle{chapterpage}
\newpage
\listoffigures
\thispagestyle{chapterpage}
\newpage
\chapter*{Resumen}
\thispagestyle{chapterpage}

La ultima vez que fui al supermercado, me compré café 100\% colombiano, o al menos eso me aseguraba la pegatina del envase... Desde entonces estuve pensando en todo el proceso que se había desarrollado con el café desde su supuesto origen hasta que llegó a mis manos. La recogida del café, su secado, el posterior envasado, hasta la venta al supermercado. Pasando por supuesto por el transporte desde el otro lado del mundo y las medidas de calidad del producto.

A todo este proceso se llama trazabilidad de un producto, el cual podemos definir como una serie de pasos para transformar materias primas o productos prefabricados en productos terminados.

Si miras a tu alrededor podrás ver una infinidad de productos que han tenido una trazabilidad hasta llegar a tus manos (y por que no seguirán teniendo más cuando termines de usarlos). Sin ir más lejos el movil o portatil que estas utilizando para leer esto tiene una trazabilidad.

Pero volviendo al café que me compré el otro día, ¿qué pasa con los usuarios? ¿Cómo sabemos que país de origen ha tenido el producto? ¿Materiales fabricados? ¿Es nuevo o usado?

Para responder a estas preguntas se ha realizado este proyecto, para dar una solución a los usuarios y que conozcan las respuestas a todas estas preguntas utilizando tecnología puntera como es Blockchain, el cual hablaremos más detenidamente en el desarrollo de este trabajo junto con la construcción de un lenguaje de dominio para la generación de una página web que utiliza la API desarrollada por Telefonica y que le han nombrado TrustOS, para poder utilizar la tecnología Blockchain sin saber nada al respecto.

\thispagestyle{empty}

\chapter*{Abstract}
\thispagestyle{empty}
EL RESUMEN EN INGLÉS


\chapter{Introducción}
\thispagestyle{chapterpage}
Esta plantilla sirve como ejemplo de TFG. Las secciones son las que están en la normativa. Esta sección la aprovechamos para introducir brevemente \LaTeX.
\par

\section{Secciones}
Las secciones se crean con \textit{section}. Las subsecciones y subsubsecciones con \textit{subsection} y \textit{subsubsection}, respectivamente. Si se desea que alguna subsección en concreto no salga en el índice se pueden usar los mismos comandos añadiéndoles un asterisco al final (\textit{subsection*} y \textit{subsubsection*}).
\par
Lorem ipsum dolor sit amet, consectetur adipiscing elit, sed eiusmod tempor incidunt ut labore et dolore magna aliqua. Ut enim ad minim veniam, quis nostrud exercitation ullamco laboris nisi ut aliquid ex ea commodi consequat. Quis aute iure reprehenderit in voluptate velit esse cillum dolore eu fugiat nulla pariatur. Excepteur sint obcaecat cupiditat non proident, sunt in culpa qui officia deserunt mollit anim id est laborum.
\subsection{una subsección}
Ut enim ad minim veniam, quis nostrud exercitation ullamco laboris nisi ut aliquid ex ea commodi consequat. Quis aute iure reprehenderit in voluptate velit esse cillum dolore eu fugiat nulla pariatur.
\subsection{otra subsección}
Excepteur sint obcaecat cupiditat non proident, sunt in culpa qui officia deserunt mollit anim id est laborum.
\subsubsection{subsubsecciones}
¡Por defecto las subsubsecciones no aparecen en el índice! Este comportamiento se puede cambiar.

\section{Citas y referencias}
Para citar un texto hay que incluirlo en el fichero \textit{LocalBibliography.bib} de esta plantilla. Una vez hecho se puede referenciar usando el comando \textit{cite}~\cite{LaTeX_tutorials}.
Para referenciar imágenes, secciones o tablas se usa el comando~\textit{ref}. Por ejemplo~\ref{fig:logoEpcc}. Es importante haber añadido una etiqueta con el nombre correspondiente al emenento a referenciar.
\par
Para facilitar la lectura, cuando se insertan citas y referencias es conveniente insertar un caracter de espacio sin salto de línea \textit{\~} antes del comando de cita o referencia.


\section{Imágenes}
Las imágenes se insertan así:

\begin{figure}[H]
    \centering
    \includegraphics[width=0.6\textwidth]{pictures/logoEpcc.jpg}
    \caption[Logo Epcc]{Logo Epcc.\\Fuente:http://www.unex.es/conoce-la-uex/centros/epcc/}
    \label{fig:logoEpcc}
\end{figure}

\chapter{Objetivos}
Lorem ipsum dolor sit amet, consectetur adipiscing elit, sed eiusmod tempor incidunt ut labore et dolore magna aliqua. Ut enim ad minim veniam, quis nostrud exercitation ullamco laboris nisi ut aliquid ex ea commodi consequat. Quis aute iure reprehenderit in voluptate velit esse cillum dolore eu fugiat nulla pariatur. Excepteur sint obcaecat cupiditat non proident, sunt in culpa qui officia deserunt mollit anim id est laborum.
\chapter{Estado del Arte}
\thispagestyle{chapterpage}
Lorem ipsum dolor sit amet, consectetur adipiscing elit, sed eiusmod tempor incidunt ut labore et dolore magna aliqua. Ut enim ad minim veniam, quis nostrud exercitation ullamco laboris nisi ut aliquid ex ea commodi consequat. Quis aute iure reprehenderit in voluptate velit esse cillum dolore eu fugiat nulla pariatur. Excepteur sint obcaecat cupiditat non proident, sunt in culpa qui officia deserunt mollit anim id est laborum.


\chapter{Metodología}
\thispagestyle{chapterpage}
Lorem ipsum dolor sit amet, consectetur adipiscing elit, sed eiusmod tempor incidunt ut labore et dolore magna aliqua. Ut enim ad minim veniam, quis nostrud exercitation ullamco laboris nisi ut aliquid ex ea commodi consequat. Quis aute iure reprehenderit in voluptate velit esse cillum dolore eu fugiat nulla pariatur. Excepteur sint obcaecat cupiditat non proident, sunt in culpa qui officia deserunt mollit anim id est laborum.
\chapter{Implementación y desarrollo}
Lorem ipsum dolor sit amet, consectetur adipiscing elit, sed eiusmod tempor incidunt ut labore et dolore magna aliqua. Ut enim ad minim veniam, quis nostrud exercitation ullamco laboris nisi ut aliquid ex ea commodi consequat. Quis aute iure reprehenderit in voluptate velit esse cillum dolore eu fugiat nulla pariatur. Excepteur sint obcaecat cupiditat non proident, sunt in culpa qui officia deserunt mollit anim id est laborum.

\chapter{Resultados}
\thispagestyle{chapterpage}
Lorem ipsum dolor sit amet, consectetur adipiscing elit, sed eiusmod tempor incidunt ut labore et dolore magna aliqua. Ut enim ad minim veniam, quis nostrud exercitation ullamco laboris nisi ut aliquid ex ea commodi consequat. Quis aute iure reprehenderit in voluptate velit esse cillum dolore eu fugiat nulla pariatur. Excepteur sint obcaecat cupiditat non proident, sunt in culpa qui officia deserunt mollit anim id est laborum.
\chapter{Conclusiones y trabajos futuros}
Lorem ipsum dolor sit amet, consectetur adipiscing elit, sed eiusmod tempor incidunt ut labore et dolore magna aliqua. Ut enim ad minim veniam, quis nostrud exercitation ullamco laboris nisi ut aliquid ex ea commodi consequat. Quis aute iure reprehenderit in voluptate velit esse cillum dolore eu fugiat nulla pariatur. Excepteur sint obcaecat cupiditat non proident, sunt in culpa qui officia deserunt mollit anim id est laborum.

\thispagestyle{empty}
\pagestyle{empty}
%%%% https://en.wikibooks.org/wiki/LaTeX/Bibliography_Management
\addcontentsline{toc}{chapter}{Bibliografía}
\bibliographystyle{unsrt}
\bibliography{LocalBibliography.bib}


\pagebreak
\appendix
\clearpage
\appendixpage
\addappheadtotoc

\chapter{Ejemplo de anexo}
Si no se desea incluir anexos, sólo hay que borrar este capítulo.
\par
Lorem ipsum dolor sit amet, consectetur adipiscing elit, sed eiusmod tempor
incidunt ut labore et dolore magna aliqua. Ut enim ad minim veniam, quis nostrud exercitation ullamco laboris nisi ut aliquid ex ea commodi consequat. Quis aute iure reprehenderit in voluptate velit esse cillum dolore eu fugiat nulla pariatur. Excepteur sint obcaecat cupiditat non proident, sunt in culpa qui officia deserunt mollit anim id est laborum.\\


\end{document}
